% SOURCE_STATEMENT: "Can the Hashin–Shtrikman bounds be improved, or can one construct microstructures that attain the Hashin–Shtrikman shear-modulus bound without attaining the bulk-modulus bound and remain optimal under finite shear? Separately, can non-macroscopically-affine deformations of periodic rod–hinge structures be obtained as limits of periodic deformations whose unit cells grow relative to the underlying structure?"
% Preparation and final self-audit: CONTRIBUTING.md prompts C and D; source/audit artifacts retained in the local physics_statements_2026_10_11 batch.
% !TeX program = lualatex
\documentclass[11pt]{article}
\usepackage[margin=25mm]{geometry}
\usepackage{fontspec}
\setmainfont{Latin Modern Roman}
\usepackage{amsmath,amssymb,mathtools}
\usepackage{unicode-math}
\setmathfont{Latin Modern Math}
\usepackage{xurl,hyperref}
\hypersetup{colorlinks=true,urlcolor=blue}
\setcounter{secnumdepth}{0}
\setlength{\parindent}{0pt}
\setlength{\parskip}{5pt}
\setlength{\emergencystretch}{3em}
\newcommand{\sref}[2]{\hyperlink{source:#1:#2}{[S#2]}}
\newcommand{\cataloguescope}[1]{\paragraph{Recorded target.}#1}
\newcommand{\R}{\mathbb R}
\newcommand{\C}{\mathbb C}
\newcommand{\N}{\mathbb N}
\newcommand{\Z}{\mathbb Z}
\newcommand{\Q}{\mathbb Q}
\newcommand{\F}{\mathbb F}
\newcommand{\eps}{\varepsilon}
\begin{document}
\section{Finite-Shear Hashin–Shtrikman Bounds and Rod–Hinge Deformations}\label{statement}

\subsection{Definitions and mathematical statement}

\paragraph{Shared notation.}$\mathbb N=\{1,2,\ldots\}$; the standard number fields have their usual meanings. All additional conventions are stated locally.


Hashin--Shtrikman bounds constrain effective moduli of composites with fixed phase data. Bulk and shear moduli measure volume and shape response, respectively. A rod--hinge framework is a periodic bar structure with movable joints. Macroscopically affine means an imposed deformation is affine at the large scale; a growing-cell approximation allows larger periodicity cells.

\paragraph{Catalogue wording retained for traceability.}
Can the Hashin–Shtrikman bounds be improved, or can one construct microstructures that attain the Hashin–Shtrikman shear-modulus bound without attaining the bulk-modulus bound and remain optimal under finite shear? Separately, can non-macroscopically-affine deformations of periodic rod–hinge structures be obtained as limits of periodic deformations whose unit cells grow relative to the underlying structure?

\paragraph{Scope and source qualification.}The two questions are separate. Attaining a shear bound without the bulk bound and remaining optimal under finite shear are both required in the first construction alternative.
\subsection{Short English statement}
Find sharper or attainable finite-shear modulus bounds, and determine whether growing periodic cells approximate nonaffine rod--hinge deformations.
\subsection{Sources}
\begin{itemize}
\item[\hypertarget{source:statement:1}{S1}] ULAM.AI and the UnsolvedMath Contributors, UnsolvedMath, catalogue entry. Catalogue curation under CC BY 4.0; definitions and English exposition added locally.
\url{https://huggingface.co/datasets/ulamai/UnsolvedMath}.
\item[\hypertarget{source:statement:2}{S2}] Mechanics of Materials: Multiscale Design of Advanced Materials and Structures (2024), Oberwolfach Reports, DOI 10.4171/owr/2024/11, pp. 657--730. Original report and mathematical context. Incorporated report material remains under the source CC BY-SA 4.0 terms unless otherwise noted.
\url{https://doi.org/10.4171/owr/2024/11}.
\end{itemize}
\label{end:statement}
\end{document}
